\documentclass[11pt]{article}
\usepackage[margin=1in]{geometry}
\usepackage{amsmath,amssymb}
\usepackage{booktabs,tabularx}
\usepackage[colorlinks=true,allcolors=blue]{hyperref}

\title{Adversarial Second-Pass Review \\
\large ``Toward an Algorithmic Free Energy Principle for Explanatory Programs''}
\author{Simulated hostile reviewer (JAIR, Round 2)}
\date{\today}

\begin{document}
\maketitle

\section*{Recommendation}

\textbf{Reject with encouragement to resubmit after empirical validation and a strengthened consistency result.} If the editor is inclined toward leniency: \textbf{Major Revision}, conditional on the authors returning with at least one empirical demonstration of the behavioral proxy and a proved (not conjectured) consistency statement in the interactive setting.

\section*{Summary of the Revision}

The authors have revised substantively. The abstract and introduction now make modest, specific claims; the related-work section is reorganized into four clusters with explicit ``Delta'' paragraphs; notation has been cleaned; a semimeasure-completion lemma has been added; posterior consistency has been correctly downgraded to a conjecture; the mechanistic-challenge section has been expanded into a concrete compliance protocol with three tests, tolerances, and a power-analysis sketch; and the tension between prediction and explanation is framed with a more careful example. The bibliography has been broadened to cover information-theoretic bounded rationality, compression-based curiosity, Bayesian experimental design, and modern interpretability tooling.

These are real improvements. The prior review's main surface-level concerns (T1--T10, the over-strong ``principled prior'' language, the aspirational mechanistic challenge) have largely been addressed. However, several structural concerns persist, and two new concerns have emerged from reading the revision closely. As a JAIR paper competing with full-theorem contributions and with empirically grounded methods papers, this submission still has a steep hill to climb.

\section*{Remaining Major Concerns}

\subsection*{RMC1. The central asymptotic result is still a conjecture.}

Section~8 (Reference-Machine Invariance) states posterior consistency in the interactive realizable case as a remark labeled ``Realizable-case consistency conjecture.'' This is honest, and an improvement over claiming a result the paper does not prove. But from a JAIR reviewer's standpoint, a theory paper whose most consequential asymptotic claim remains unproved is a theory paper that is not finished. The variational identity (Theorem~3.2) and the MDL/MAP bridge (Section~5) are correct, but they are essentially bookkeeping: restatements, under careful notation, of standard facts about Gibbs posteriors under $2^{-|p|}$ priors. The one genuinely new quantitative claim the framework needs is the interactive-setting consistency, and it is absent.

\emph{What would fix this}: Either (a) prove the conjecture under explicit exploration and regularity conditions (extending Blackwell--Dubins and Leike's passive arguments to interactive semimeasures is a real technical step, but not obviously out of reach); or (b) prove a weaker but non-trivial regret or posterior-concentration bound at finite $t$ whose constants the paper can bear.

\subsection*{RMC2. No empirical content (still).}

The paper explicitly defers empirical validation to a companion paper. For a pure-theory submission this would be acceptable, but the paper does not read as a pure-theory submission: it invokes ARC, DreamCoder, and FunSearch; it proposes a behavioral proxy; it specifies a compliance protocol with tolerances; it promises that the framework ``makes prediction and explanation mathematically distinguishable.'' All of these are empirical posture without empirical delivery. At JAIR, where methods papers are routinely evaluated against stronger-than-theory baselines, this gap is the single most important determinant of the review outcome.

The right reading is not ``the theorems are wrong'' but ``the theorems on their own do not justify the venue.'' JAIR readers will want at least a toy demonstration in which an agent that optimizes the behavioral proxy provably outperforms a prediction-only baseline on held-out inputs that are structurally different from the query set.

\subsection*{RMC3. The synthesis-versus-contribution question has not been fully retired.}

The revision acknowledges that ``none of these steps is individually unprecedented'' and positions the paper as a specific assembly. This is the correct and honest framing. But it concedes the most damaging reviewer argument in the prior round: that the contribution is incremental combination rather than new ground. In philosophy of AI, Ortega--Braun already have the variational bounded-rationality structure; Schmidhuber has compression-based curiosity; Lindley--Chaloner--Verdinelli--MacKay have the experimental-design reading; Hutter and Leike have the universal-prior, interactive-semimeasure apparatus. The paper argues that the specific combination---Solomonoff prior inside a variational functional with a Lindley-style action rule---is new. This is true, but the reviewer's question is whether that specific combination produces a new \emph{result} the individual components do not. Corollary~6.6 (risk-minus-IG under deterministic discovery) recovers Lindley; Theorem~3.2 recovers the standard Gibbs identity; the MDL/MAP bridge restates the standard MAP-as-MDL fact. What new quantity or prediction does the synthesis yield that an Ortega--Braun practitioner would not already have access to?

The honest answer, I think, is: the three-level framework itself (ideal / behavioral proxy / mechanistic challenge). But that is a conceptual contribution, not a technical one, and it rides on the protocol in Section~8 actually being run.

\subsection*{RMC4. The mechanistic challenge remains unexercised.}

Section~8 is now careful and specific: three quantities, three verification tests, explicit tolerances, a statistical-power sketch, and an interpretability-substrate subsection connecting to real tooling (linear probes, SAEs, ROME). This is well done. But the section reads as a specification of a protocol rather than a run of one. Reviewers will ask the obvious question: does the protocol, when applied to at least one small model (a trained DSL interpreter, a DreamCoder-style agent, a TinyTransformer), actually discriminate compliance from non-compliance? A compliance protocol that has never been applied to a real system is a compliance protocol whose calibration is unknown.

\subsection*{RMC5. The ``hard-to-vary'' criterion is under-developed relative to the weight placed on it.}

Deutsch's hard-to-vary idea is invoked in Section~2, promoted to one of the three mechanistic-challenge tests in Section~8, and returned to in the Discussion. But the paper's formal apparatus attaches compactness (Solomonoff prior) to hidden programs and reuse (held-out evaluation) to the proxy, while hard-to-vary structure is offloaded to the mechanistic challenge as a decay-rate property under DSL edit distance. A reviewer who cares about this criterion will note that (i) DSL edit distance is itself a modeling choice, (ii) the exponential-decay threshold $\alpha$ is unmotivated except as a research parameter, and (iii) there is no theorem connecting the Solomonoff prior to hard-to-vary structure. One would expect a lemma of the form: under the universal prior, typical programs near a MAP explanation have substantially worse posterior mass, with rate $f(|\Delta|)$. Without such a link, hard-to-vary remains a fourth criterion that the framework gestures at but does not operationalize formally.

\section*{New Concerns Emerging from the Revision}

\subsection*{N1. Rhetorical slippage between ``variationally specified'' and ``normatively complete.''}

The revision replaces ``normatively complete'' with ``variationally specified.'' This is a defensible move, but in several places the paper's prose still carries the earlier weight---for example, the abstract's claim that the theory ``combines two traditions that have remained partially separate'' and the conclusion's claim that the result is ``a variational and active theory in which hidden causes are explanatory programs and prior mass over those programs is induced by self-delimiting description length.'' These are factually true descriptions of the formal object defined, but they still suggest the synthesis has settled a question rather than organized an open research program. A hostile reviewer will hold the authors to the ``variationally specified'' standard throughout and will flag every sentence that slides toward normativity.

\subsection*{N2. The reference-machine invariance argument partly concedes its own motivation.}

Section~8 is admirably frank: the invariance constant $c_{U,V}$ ``can be arbitrarily large on any fixed finite-sample problem.'' This is correct. But the paper's main motivational contrast is between a ``stipulated FEP prior'' and a ``principled'' universal prior. If the universality constant is unbounded at any finite sample, then on any finite experiment the choice of universal machine can shift the induced prior by an arbitrary amount---which is exactly the kind of domain-specific stipulation the paper wanted to avoid. The paper partly addresses this in the Discussion, noting that the claim of improvement is ``one can reasonably accept or question,'' but a reviewer may argue that this honest caveat undermines the framing rather than qualifying it. At minimum, the abstract and conclusion should be aligned with this caveat; right now they are not.

\subsection*{N3. Some cited recent work is load-bearing but the citations are thin.}

The paper leans on very recent interpretability work (Cunningham et al.\ 2023 on SAEs, Templeton et al.\ 2024 on scaling, Olsson et al.\ 2022 on induction heads, Wang et al.\ 2023 on circuits, Meng et al.\ 2022 on ROME) to make the claim that the mechanistic challenge is ``not in principle intractable.'' The citations are appropriate but the claims are load-bearing and deserve more than one-sentence treatment. A short paragraph explaining what SAEs or ROME actually expose, and how that relates concretely to exposing $q_t$ or a program neighborhood, would make the argument defensible.

\subsection*{N4. The Ortega--Braun comparison is weak in one specific respect.}

The Delta paragraph for the information-theoretic bounded rationality cluster says: ``Ortega--Braun give the variational form of bounded-rational control but work with generic priors; we specialize the prior to the Solomonoff universal measure over interactive programs.'' This is true but incomplete. Ortega--Braun's KL-regularized control is derived from thermodynamic/free-energy principles with specific information-processing cost interpretations. Our $\tau$ parameter plays a structurally similar role. The paper should either (a) say more precisely how our scaled functional $\Fcal_t^\tau$ relates to the Ortega--Braun free-energy-as-KL-control objective (do they coincide under a specific identification?), or (b) acknowledge that this is a detailed open question and point to where an interested reader could take it up. Right now the paper treats Ortega--Braun as a component whose technical form is orthogonal to ours, which is likely too clean a picture.

\section*{Where the Revision Clearly Succeeds}

To keep the record balanced:

\begin{itemize}
\item The related-work clustering with ``Delta'' paragraphs is the right architecture for a positioning-heavy theory paper.
\item The Running Example is concrete, non-linear, and has the branch/modular structure that makes prediction-by-lookup visibly insufficient.
\item The semimeasure-completion lemma is a clean response to a real technical concern about mutual-information identities under subnormalized kernels.
\item The three-level framework (ideal / proxy / mechanistic) is a useful organizing move that other papers in this space lack.
\item The mechanistic-challenge compliance protocol is the single most valuable contribution of the paper and is the clearest advance over the prior round.
\item The naive-future-free-energy proposition (showing predictive entropy minimization is not the right epistemic objective) is clean and pedagogically useful.
\end{itemize}

\section*{Verdict}

As a self-contained theoretical contribution with (i) Theorem~3.2, (ii) the MDL/MAP bridge, (iii) the risk-minus-IG decomposition, and (iv) a specified mechanistic-challenge protocol, the paper has a coherent core. What it lacks, by the standards of a top AI journal, is a load-bearing new technical result (RMC1), empirical demonstration (RMC2), and validation of its compliance protocol (RMC4). If the editor weights conceptual clarity and positioning highly, this paper is a Major-Revision-to-Acceptance candidate. If the editor weights novel technical results or empirical teeth, this paper is a Reject with encouragement to resubmit.

\end{document}
